% Part:sets-functions-relations
% Chapter: size-of-sets
% Section: non-enumerability-alt

\documentclass[../../../include/open-logic-section]{subfiles}

\begin{document}

\olfileid{sfr}{siz}{nen-alt}

\olsection{\printtoken{S}{nonenumerable} سیٹ}

\begin{editorial}
  ایہ حصہ \olref[enm-alt]{sec} دیاں تعریفاں ورت کے $\Bin^\omega$ تے
  $\Pow{\Nat}$ دے گنتی وار فہرست نہ بنن جوگ ہون نوں ثابت کردا اے،
  یعنی~$\Nat$ نال اک بہ اک مطابقت دی لوڑ رکھ کے،~$\PosInt$ توں ہر
  ہدف قدر لین والے نقشے دی بجائے۔
\end{editorial}

\begin{explain}
قدرتی عدداں دا سیٹ~$\Nat$ لامتناہی اے۔ ایہ صاف طور تے
!!{enumerable} وی اے۔ پر حیران کرن والی حقیقت ایہ اے کہ
\emph{!!{nonenumerable}} سیٹ وی نیں، یعنی ایہو جہے سیٹ جیہڑے
!!{enumerable} نہیں (ویکھو \olref[sfr][siz][enm-alt]{defn:enumerable})۔

ایہ گل حیران کر سکدی اے۔ آخرکار، ایہ آکھنا کہ~$A$،
!!{nonenumerable} اے، ایہ آکھنا اے کہ کوئی !!{bijection}
$f \colon \Nat \to A$ \emph{موجود نہیں}؛ یعنی~$\Nat$ دے لامتناہی
!!{element}s نوں~$A$ ول نقشہ بند کرن والا کوئی فنکشن سارے~$A$ نوں
پورا نہیں کردا۔ سو جے~$A$، !!{nonenumerable} اے، تاں~$A$ دے
!!{element}s قدرتی عدداں نالوں ”ودھ“ نیں۔

کسی سیٹ دے !!{nonenumerable} ہون نوں ثابت کرن لئی تہانوں وکھاؤنا
پئے گا کہ کوئی مناسب !!{bijection} موجود نہیں ہو سکدی۔ ایہ کرن دا
سب توں چنگا طریقہ ایہ وکھاؤنا اے کہ~$A$ دے !!{element}s دی گنتی وار
فہرست بناؤن دی ہر کوشش گھٹ توں گھٹ اک !!{element} ضرور چھڈ دیندی
اے؛ ایہ وکھاؤندا اے کہ کوئی فنکشن $f\colon \Nat \to A$،
!!{surjective} نہیں۔ اینوں قائم کرن دی اک عام حکمتِ عملی کانتور دا
\emph{قطری طریقہ} ورتنا اے۔~$A$ دے !!{element}s دی کوئی فہرست،
مثلاً $x_1$، $x_2$، \dots، دتی ہووے، تاں اسی~$A$ دا اک ہور
!!{element} بناؤندے آں جیہڑا اپنی بناوٹ کر کے اوس فہرست وچ ہو ای
نہیں سکدا۔

پر ایہ ساری گل مثال نال سب توں چنگی طرح سمجھ آندی اے۔ سو ساڈی پہلی
مثال سیٹ~$\Bin^\omega$ اے، جیہڑا $0$ تے $1$ دیاں ساریاں لامتناہی
لڑیاں دا سیٹ اے۔ ($\Bin$ توں مراد ثنائی اے، تے اسی اینوں بس دو
دو رکناں والے سیٹ
$\{0,1\}$ دے طور تے سمجھ سکدے آں۔)\oliflabeldef{sfr:card-arithmetic:card-opps:sec}{\footnote{ہور
ٹھیک طور تے، سانوں ایہ شرط لانی چاہیدی اے کہ~$\Bin^\omega$،
$\omega$-لڑیاں دا اوہ سیٹ اے جیہڑیاں $0$ تے $1$ دیاں نیں، یعنی سیٹ
$\funfromto{\omega}{\{0,1\}}$۔ پر ایس دا مطلب
\olref[sfr][card-arithmetic][card-opps]{sec} وچ جا کے صاف ہووے گا۔}}{}
پر فی الحال ایہ ذرا ڈھیلی عبارت کوئی بھلیکھا نہیں پاؤنی چاہیدی۔
\end{explain}

\begin{thm}
\ollabel{thm:nonenum-bin-omega}
$\Bin^\omega$، !!{nonenumerable} اے۔
\end{thm}

\begin{proof}
$\Bin^\omega$ دے کسے ذیلی سیٹ دی کوئی وی گنتی وار فہرست لو۔ سو ساڈے
کول کوئی فہرست $s_{0}$، $s_{1}$، $s_{2}$، \dots{} اے، جتھے ہر
$s_n$، $0$ تے~$1$ دی اک لامتناہی لڑی اے۔ آؤ $s_n(m)$ نوں ایس فہرست
دی $n$ویں لڑی دا $m$واں ہندسہ آکھیے۔ ہُن اسی اپنی فہرست نوں اک جدول
سمجھ سکدے آں، جتھے $s_n(m)$ نوں $n$ویں قطار تے $m$ویں کالم وچ رکھیا
گیا اے:
% OLSIZ-008 ماخذ دی درستی: انگریزی نثر اشارییاں دے کردار اُلٹ دَسدی اے؛ اگلا جملہ تے جدول s_n(m) نوں nویں لڑی دا mواں ہندسہ مقرر کردے نیں۔
\[
\begin{array}{c|c|c|c|c|c}
& 0 & 1 & 2 & 3 & \dots \\\hline
0 & \mathbf{s_{0}(0)} & s_{0}(1) & s_{0}(2) & s_0(3) & \dots \\\hline
1 & s_{1}(0)& \mathbf{s_{1}(1)} & s_1(2) & s_1(3) & \dots \\\hline
2 & s_{2}(0)& s_{2}(1) & \mathbf{s_2(2)} & s_2(3) & \dots \\\hline
3 & s_{3}(0)& s_{3}(1) & s_3(2) & \mathbf{s_3(3)} & \dots \\\hline
\vdots & \vdots & \vdots & \vdots & \vdots & \mathbf{\ddots}
\end{array}
\]
ہُن اسی اک لامتناہی لڑی~$d$، جیہڑی $0$ تے $1$ دی اے، بناواں گے
جیہڑی ایس فہرست وچ نہیں۔ اسی اوہدی ہر قدر دَس کے ایہ کم کراں گے،
یعنی اسی $d(n)$ مقرر کراں گے، ہر $n \in \Nat$ لئی۔ بدیہی طور تے، اسی اُتلے جدول
دا قطر تھلے ول پڑھدے آں (ایسے لئی ناں ”قطری طریقہ“ اے)، تے فیر ہر
$1$ نوں $0$ تے ہر $0$ نوں~$1$ وچ بدل دیندے آں۔
% OLSIZ-009 ماخذ دی درستی: انگریزی نثر 1 توں 0 والی تبدیلی دو واری آکھدی اے؛ اگلی حالت وار تعریف موجب دوجی تبدیلی 0 توں 1 اے۔
ہور مجرد طور تے، اسی $d(n)$ نوں $0$ یا $1$ ایس مطابق تعریف کردے آں
کہ قطر دا $n$واں !!{element}، $s_n(n)$، $1$ اے یا $0$، یعنی:
\[
d(n) =
\begin{cases}
1 & \text{جے $s_{n}(n) = 0$ اے}\\
0 & \text{جے $s_{n}(n) = 1$ اے}
\end{cases}
\]
صاف اے کہ $d \in \Bin^\omega$، کیوں جے ایہ $0$ تے $1$ دی اک
لامتناہی لڑی اے۔ پر اسی~$d$ نوں ایویں بنایا اے کہ
$d(n) \neq s_n(n)$، ہر $n \in \Nat$ لئی۔ یعنی~$d$، $s_n$ نال
اوہدی $n$ویں قدر وچوں وکھ اے۔ سو $d \neq s_n$، ہر
$n\in \Nat$ لئی۔ ایس لئی
$d$ فہرست $s_0$، $s_1$، $s_2$،
\dots وچ نہیں ہو سکدا۔

اسی وکھا دتا اے کہ~$\Bin^\omega$ دے کسے ذیلی سیٹ دی من مانی گنتی
وار فہرست دتی ہووے، تاں اوہ~$\Bin^\omega$ دا کوئی !!{element} چھڈ
دیوے گی۔ سو سیٹ~$\Bin^\omega$ دی کوئی گنتی وار فہرست نہیں، یعنی
$\Bin^\omega$، !!{nonenumerable} اے۔
\end{proof}

\begin{explain}
ایس ثبوت دے طریقے نوں ”قطری بناوٹ“ آکھیا جاندا اے، کیوں جے ایہ~$d$
نوں تعریف کرن لئی جدول دا قطر ورتدا اے۔ پر قطری بناوٹ وچ جدول دا
موجود ہونا لازمی نہیں۔ اصل وچ، اسی اوسے ورگا خیال ورت کے کسی سیٹ دا
!!{nonenumerable} ہونا وی وکھا سکدے آں، بھاویں کوئی جدول تے حقیقی قطر
شامل نہ ہووے۔ اگلا نتیجہ دَسدا اے کہ ایہ کِویں ہوندا اے۔
\end{explain}

\begin{thm}
\ollabel{thm:nonenum-pownat}
$\Pow{\Nat}$، !!{enumerable} نہیں۔
\end{thm}

\begin{proof}
اسی اوسے ڈھنگ نال کم کردے آں، ایہ وکھا کے کہ~$\Nat$ دے ذیلی سیٹاں دی
ہر فہرست~$\Nat$ دا کوئی ذیلی سیٹ چھڈ دیندی اے۔ سو فرض کرو ساڈے کول
کوئی فہرست $N_0, N_1, N_2, \ldots$ اے، جیہڑی~$\Nat$ دے ذیلی سیٹاں
دی اے۔ اسی
سیٹ~$D$ ایویں تعریف کردے آں: $n \in D$ عین اوس ویلے جدوں
$n \notin N_{n}$:
\[
D = \Setabs{n \in \Nat}{n \notin N_n}
\]
صاف اے کہ $D\subseteq \Nat$۔ پر~$D$ فہرست وچ نہیں ہو سکدا۔ آخرکار،
بناوٹ موجب $n \in D$ عین اوس ویلے جدوں $n\notin N_n$، ایس لئی
$D \neq N_n$، ہر $n \in \Nat$ واسطے۔
\end{proof}

\begin{explain}
پچھلے ثبوت وچ قطر دا ذکر نہیں آیا۔ فیر وی تسیں ایس نوں قطری دلیل
سمجھ سکدے او جے ایس دی تصویر ایویں بناؤ: تصور کرو کہ سیٹ $N_0$،
$N_1$، \dots، اک جدول وچ لکھے نیں، جتھے~$N_n$ نوں $n$ویں قطار وچ
لکھن لئی $m$ نوں $m$ویں کالم وچ عین اوس ویلے لکھدے آں جدوں
$m \in N_n$۔ مثال لئی، من لو کہ ایس فہرست دے پہلے چار سیٹ
$\{0,1,2,\dots\}$، $\{1, 3, 5, \dots\}$، $\{0,1,4\}$، تے
$\{2,3,4,\dots\}$ نیں؛ تاں ساڈا جدول ایویں شروع ہووے گا:
\[
\begin{array}{r@{}rrrrrrr}
  N_0 = \{ & \mathbf{0}, & 1, & 2, & & & & \dots\}\\
  N_1 = \{ &  & \mathbf{1}, &  & 3, &  & 5, & \dots\}\\
  N_2 = \{ & 0, & 1, &  &  & 4\phantom{,} &  & \}\\
  N_3 = \{ &  &  & 2, & \mathbf{3}, & 4, & & \dots\}\\
  &\vdots & & & & & & \ddots\phantom{\}}
  \end{array}
\]
فیر~$D$ اوہ سیٹ اے جیہڑا قطر دے نال تھلے جا کے، $n \in D$ عین اوس
ویلے رکھ کے ملدا اے جدوں~$n$ قطر اُتے \emph{نہیں}۔ سو اُتلی
صورت وچ اسی $0$ تے $1$ چھڈ دیاں گے،~$2$ شامل کراں گے،~$3$ چھڈ
دیاں گے، تے ایویں اگے۔
\end{explain}

\begin{prob}
کھلی قطری دلیل نال وکھاؤ کہ سارے فنکشناں $f \colon \Nat \to \Nat$
دا سیٹ !!{nonenumerable} اے۔ یعنی وکھاؤ کہ جے $f_1$، $f_2$، \dots،
فنکشناں دی فہرست اے تے ہر $f_i\colon \Nat \to \Nat$، تاں کوئی
$g \colon \Nat \to \Nat$ ایہو جہا اے جیہڑا ایس فہرست وچ نہیں۔
\end{prob}

\end{document}
